\chapter{多分类学习的可学习性}
\label{chap:multiclass-learnability}

二分类只有两种答案。一个假设在某个输入上改变预测时，改变后的标签已经确定；
多分类则允许它改成多个不同标签。这一差别影响了假设空间的计数，也改变了可学习性
与一致收敛之间的关系。本章先讨论标签数有限时如何推广VC维，再说明无限标签为什么
需要区分图维与DS维，最后考察这些统计结论在要求学习程序可计算时需要哪些补充条件。

\section{从二分类到多分类}
\label{sec:multiclass-setting}

以邮件分类为例，标签可以是正常邮件、广告或诈骗。输入$x$仍然表示邮件特征，
假设$h$为它选择一个标签；预测标签与真实标签不同，就记一次错误。
设输入空间为$X$、标签空间为$Y$，假设空间为非空集合$\mathcal H\subseteq Y^X$。
对输入与标签的联合分布$\mathcal P$以及独立同分布样本$S=((x_i,y_i))_{i=1}^n$，
零一损失下的期望风险与经验风险为
\begin{equation}
 R_{\mathcal P}(h)=\mathbb P_{(x,y)\sim\mathcal P}\{h(x)\neq y\},\qquad
 \hat R_S(h)=\frac1n\sum_{i=1}^n\mathbb I\{h(x_i)\neq y_i\}.
 \label{eq:multiclass-risk}
\end{equation}
本章的风险均使用零一损失。PAC和不可知PAC的成功标准沿用前章：前者在可实现分布上
要求期望风险足够小，后者要求期望风险接近$\mathcal H$所能达到的最小值。
以下统计定理均假定所需的可测性条件成立。

有限个类别确实可以用多个二分类标签的组合来表示。例如，将三类邮件编码为
\[
 \text{正常}\longleftrightarrow(0,1),\qquad
 \text{广告}\longleftrightarrow(1,0),\qquad
 \text{诈骗}\longleftrightarrow(1,1).
\]
这样，一封标为“广告”的邮件就为两个二分类问题分别提供标签1和0。我们用同一批邮件
训练两个二分类器$b_1,b_2$；面对新邮件$x$，先得到$(b_1(x),b_2(x))$，再按编码表
还原类别。若得到未使用的编码$(0,0)$，可以约定输出一个固定类别，例如正常邮件。
这里把一个三分类问题转成了两个二分类问题，并没有减少需要辨认的类别信息。

这种做法可以推广到任意有限标签空间。若$|Y|=K$，用
$L=\lceil\log_2K\rceil$个二进制位就足以为所有类别分配不同编码。
记编码为$c:Y\to\{0,1\}^L$，第$j$位为$c_j(y)$。
原多分类假设空间$\mathcal H$在这一位上诱导的二分类假设空间是
\[
 \mathcal B_j=\{x\mapsto c_j(h(x)):h\in\mathcal H\},\qquad 1\leq j\leq L.
\]
这个式子说明二分类子问题的假设从哪里来：把每个多分类假设的输出标签换成编码，
再只保留第$j$位。标签虽然变成了0和1，$\mathcal B_j$本身仍可能很复杂；
二分类理论要求它具有有限VC维，才保证分布无关的PAC可学习性。
编码本身不会自动满足这个要求。

在可实现情形，如果每个$\mathcal B_j$都PAC可学习，这一归约就能给出多分类保证。
设学到的第$j$个二分类器的期望风险为
\[
 e_j=\mathbb P_{(x,y)\sim\mathcal P}\{b_j(x)\neq c_j(y)\}.
\]
令$\tilde h$表示编码还原后的分类器。当所有位都预测正确时，$\tilde h$必然输出真实
类别。因此，多分类预测错误只能发生在至少一位预测错误时，从而
\[
 R_{\mathcal P}(\tilde h)\leq\sum_{j=1}^L e_j.
\]
把各二分类器的期望风险分别控制在$\varepsilon/L$以内，各自学习失败的概率分别
控制在$\delta/L$以内，就能以至少$1-\delta$的概率保证
$R_{\mathcal P}(\tilde h)\leq\varepsilon$。各二分类器可以使用同一批样本，
这里不要求它们的错误相互独立；合成后的分类器也未必属于原假设空间。

不可知情形则不能仅凭“每一位都学到最优”得到原多分类问题的最优期望风险。
仍用上面的编码，假设输入只有一个值，三类标签出现的概率依次为$0.40,0.35,0.25$，
并允许三个常量多分类假设。第一位取1的概率为$0.60$，第二位取1的概率为$0.65$，
所以两个最优二分类器都预测1，合成后始终预测诈骗，期望风险为$0.75$。
但始终预测正常的多分类假设具有更小的期望风险$0.60$。
这里每一位的最优预测已经知道，即使样本无限多，逐位优化仍然选错了最优类别。
原因是逐位学习分别比较各位的错误比例，而多分类零一损失比较的是整个类别是否正确。

因此，二进制编码提供了有限标签下的一种算法构造方法；要建立学习保证，还需分析
编码诱导的二分类假设空间，并把各位的期望风险转换为多分类期望风险。
标签无限时，任何固定长度的二进制编码都无法区分所有类别，上面的有限次归约也就
不能直接使用。为了统一讨论有限和无限标签，并直接判断原假设空间的可学习性，
我们还需要研究多分类假设在一组输入上能够实现哪些预测。

直接分析多分类假设空间时，也不能简单要求它在$m$个输入上实现全部$|Y|^m$种预测。
例如，一个假设空间可能只使用每个输入处的两个标签，却已经包含任意大的二分类问题。
若要求使用全部标签才算打散，就会漏掉这种困难。合理的推广应允许各输入只在两个
预先指定的标签之间切换，由此引出Natarajan维。

\section{有限标签空间}
\label{sec:multiclass-dimensions}
\label{sec:finite-label-learning}

本节固定$2\leq K=|Y|<\infty$。先找出多分类中的二元选择结构，再用它限制有限
输入上的预测数目，就能沿用前章从一致收敛到ERM泛化的证明。

\subsection{Natarajan维}

对两封固定邮件$x_1,x_2$，在第一封上指定正常与广告，在第二封上指定正常与诈骗。
若假设空间能给出以下四种预测，就能在这两对标签之间独立选择。
\begin{center}
\begin{tabular}{cc}
\toprule
$x_1$上的预测 & $x_2$上的预测\\
\midrule
正常 & 正常\\
广告 & 正常\\
正常 & 诈骗\\
广告 & 诈骗\\
\bottomrule
\end{tabular}
\end{center}
每个输入的两个标签可以不同，但指定之后不能随选择结果改变。

\begin{definition}[Natarajan打散与Natarajan维]
\label{def:natarajan-dimension}
设$C=\{x_1,\ldots,x_m\}\subseteq X$。若存在$f,g:C\to Y$，满足
$f(x)\neq g(x)$对所有$x\in C$成立，且对每个$B\subseteq C$都存在$h_B\in\mathcal H$使
\begin{equation}
 h_B(x)=\begin{cases}f(x),&x\in B,\\g(x),&x\in C\setminus B,\end{cases}
 \label{eq:natarajan-shattering}
\end{equation}
则称$C$被$\mathcal H$ Natarajan打散。令$\mathcal C_N$为所有被
$\mathcal H$ Natarajan打散的有限输入集组成的集合，定义
\[
 d_N(\mathcal H)=\sup_{C\in\mathcal C_N}|C|.
\]
\end{definition}

这个定义也适用于无限标签空间。当$|Y|=2$时，每个输入上的不同标签对只能是那两个
二分类标签，所以Natarajan维就是VC维。维度为零表示每个输入上的预测都唯一；
可打散的有限点集大小没有上界时，维度为无穷。

可以用一个可直接核对的假设空间理解“维度恰好为2”。设输入空间只有三封固定邮件
$x_1,x_2,x_3$，假设空间也只有图~\ref{fig:natarajan-concrete}中的四个假设。
对$x_1$选定“正常、广告”，对$x_2$选定“正常、诈骗”，四个假设恰好实现两对标签的
全部四种组合，因此$\{x_1,x_2\}$被Natarajan打散。这一步给出$d_N\geq2$。

\begin{figure}[htbp]
 \centering
 % Four complete predictions on three fixed inputs; Natarajan dimension two.
\begin{tikzpicture}[
 box/.style={draw=BookRule,fill=BookPaper,rounded corners=2pt,
 minimum width=39mm,minimum height=12mm,align=center,font=\small},
 every node/.style={font=\small}]
 \node[box] (a) at (0,0) {$h_1$：正常，正常，正常};
 \node[box] (b) at (5.7,0) {$h_2$：广告，正常，正常};
 \node[box] (c) at (0,-2) {$h_3$：正常，诈骗，正常};
 \node[box] (d) at (5.7,-2) {$h_4$：广告，诈骗，正常};
 \draw[BookBlue,thick] (a)--node[above,align=center]{只改$x_1$}(b);
 \draw[BookBlue,thick] (c)--node[below,align=center]{只改$x_1$}(d);
 \draw[BookTeal,thick] (a)--node[left]{只改$x_2$}(c);
 \draw[BookTeal,thick] (b)--node[right]{只改$x_2$}(d);
 \node[align=center,text=BookMuted] at (2.85,-3)
 {每行依次给出$x_1,x_2,x_3$的预测；$x_3$始终为正常。};
\end{tikzpicture}

 \caption{Natarajan打散的四种预测。沿横向只改变第一封邮件的预测，沿纵向只改变第二封
 邮件的预测；第三封邮件始终预测为正常。每个输入的两个标签预先固定，四个顶点实现全部组合。}
 \label{fig:natarajan-concrete}
\end{figure}
\FloatBarrier

第三封邮件在所有假设下都预测为正常，所以任何包含$x_3$的点集都无法在该输入上
实现两个不同标签。输入空间只有这三个点，不存在能被打散的三点集，因此$d_N\leq2$。
上下两步合起来才说明$d_N=2$。这与二分类VC维例子的证明方式一致：既要找到能被打散的
点集，也要证明更大的点集不能被打散。不同的是，这里两个输入使用的标签对并不相同。

\subsection{有限标签的一致收敛}

固定$m$个输入后，不同假设可能给出相同的预测向量。记
\[
 V(x_1,\ldots,x_m)
 =\{(h(x_1),\ldots,h(x_m)):h\in\mathcal H\}.
\]
生长函数取所有输入选择中最大的预测数目，即
\[
 \Pi_{\mathcal H}(m)=\sup_{(x_1,\ldots,x_m)\in X^m}|V(x_1,\ldots,x_m)|.
\]
这与第~\ref{sec:binary-growth-sauer}节的二分类生长函数是同一个计数问题：固定输入，
合并给出相同预测的假设，再对输入的选择取最大值。区别只是预测向量的每个坐标现在
可以取$K$个标签。二分类的Sauer--Shelah引理用VC维限制预测数目；下面的定理用
Natarajan维完成对应的多分类计数。

\begin{theorem}[有限标签的预测数目上界]
\label{thm:multiclass-growth-bound}
设标签数为$K<\infty$，Natarajan维为$d<\infty$。
记$q=K(K-1)/2$为不同标签对的数目，则
\[
 \Pi_{\mathcal H}(m)\leq G_{d,K}(m)
 :=\sum_{j=0}^{\min\{d,m\}}\binom mj q^j.
\]
\end{theorem}

\begin{proof}
对任意有限预测集合$V\subseteq Y^m$，删去最后一个坐标后得到前缀集合$U$。
若某个前缀能够接上$r$种末尾标签，它对$|V|$的贡献为$r$，对$|U|$的贡献为1。
把多出的$r-1$次计数分配给标签对，因为$r-1\leq r(r-1)/2$，得到
\[
 |V|\leq|U|+\sum_{\{a,b\}\subseteq Y}|U_{a,b}|,
\]
其中$U_{a,b}$表示既能接上$a$又能接上$b$的前缀集合。
$U$的Natarajan维至多为$d$；每个$U_{a,b}$的Natarajan维至多为$d-1$，
否则把最后一个坐标上的标签对$(a,b)$加入，就会在$V$中打散$d+1$个坐标。

因此最大预测数满足递推上界
\[
 M(m,d)\leq M(m-1,d)+qM(m-1,d-1).
\]
零个坐标只有一个空预测，维度为零时也至多有一个预测向量。
对$m$归纳，并使用二项式系数的递推关系，即得所列求和式。
\end{proof}

当$K=2$时，$q=1$且$d_N=\operatorname{VC}(\mathcal H)$，定理中的求和式恰好变为
二分类Sauer--Shelah引理的$\sum_{j=0}^{\min\{d,m\}}\binom mj$。
多分类中额外的$q^j$反映了每个参与二元选择的坐标还可以选择不同的标签对。

对固定$d,K$，这个上界随$m$至多按多项式增长。由于样本标签固定后，相同的预测向量
必然产生相同的零一损失向量，因此这个计数也控制经验风险的可能取值。
前章的副样本和随机交换分析便可用于多分类。

% 避免短定理在页末仅留下标题和半句条件。
\begingroup
\tcbset{lines before break=4}
\begin{theorem}[有限标签的一致收敛保证]
\label{thm:finite-label-uniform-convergence}
在前一定理的条件下，对任意$\delta\in(0,1)$和$S\sim\mathcal P^n$，
以至少$1-\delta$的概率，
\begin{equation}
 \sup_{h\in\mathcal H}|R_{\mathcal P}(h)-\hat R_S(h)|
 \leq\sqrt{\frac8n\log\frac{4G_{d,K}(2n)}{\delta}}.
 \label{eq:finite-label-uniform-bound}
\end{equation}
\end{theorem}
\endgroup

\begin{proof}
令$\eta>0$且$n\eta^2\geq2$。引入独立副样本后，对称化把期望风险与经验风险的
偏差事件，转化为两份经验风险相差超过$\eta/2$的事件，概率至多增加一倍。
在合并的$2n$个输入上，不同损失向量至多有$G_{d,K}(2n)$种。
随机交换每对观测，对每个固定向量使用Hoeffding不等式，再把这些概率相加，得到
\[
 \mathbb P\!\left\{\sup_h|R_{\mathcal P}(h)-\hat R_S(h)|>\eta\right\}
 \leq4G_{d,K}(2n)e^{-n\eta^2/8}.
\]
令右侧为$\delta$便得到结论；所得$\eta$自动满足$n\eta^2\geq2$。
若上界大于1，则由零一损失取值范围直接成立。
\end{proof}

这里的分布是$X\times Y$上的联合分布，并未要求类别出现概率相同。
固定$d,K,\delta$后，式~\eqref{eq:finite-label-uniform-bound}趋于零，
且同一个样本量适用于所有允许分布，这正是分布无关的一致收敛。
例如，$1\leq d\leq2n$时可用
$G_{d,K}(2n)\leq q^d(2en/d)^d$进一步估计右侧，看到类别数通过$\log K$影响这个上界。

\subsection{有限标签的可学习性}

一致收敛保证每个假设的经验风险都接近其期望风险，因此ERM在样本上的选择也能得到
期望风险保证。结合Natarajan打散所给出的不可学习性下界，可以得到以下等价关系。

\begin{theorem}[有限标签的可学习性刻画]
\label{thm:finite-label-learnability}
若$2\leq|Y|=K<\infty$，则以下命题等价：
\begin{enumerate}
 \item $d_N(\mathcal H)<\infty$；
 \item 零一损失满足分布无关的一致收敛；
 \item $\mathcal H$不可知PAC可学习；
 \item $\mathcal H$可实现PAC可学习。
\end{enumerate}

设$d=d_N(\mathcal H)<\infty$。对任意$\varepsilon,\delta\in(0,1)$，若
\[
 n\geq\frac{32}{\varepsilon^2}\log\frac{4G_{d,K}(2n)}{\delta},
\]
则任意ERM以至少$1-\delta$的概率满足
\[
 R_{\mathcal P}(A(S))\leq\inf_{h\in\mathcal H}R_{\mathcal P}(h)+\varepsilon.
\]
\end{theorem}

\begin{proof}
有限Natarajan维推出前一定理的一致收敛。令其精度为$\varepsilon/2$，对任意比较
假设$h$，有
\begin{align*}
 R_{\mathcal P}(A(S))&\leq\hat R_S(A(S))+\varepsilon/2\\
 &\leq\hat R_S(h)+\varepsilon/2\\
 &\leq R_{\mathcal P}(h)+\varepsilon.
\end{align*}
对$h$取下确界得到不可知PAC保证，而它当然涵盖可实现分布。

反过来，若Natarajan维无限，对任意$n$都能找到$2n$个被打散的输入，
并在每个输入处固定两个不同标签。令输入均匀分布，目标在每对标签中独立等概率选择。
$n$次抽样至多看到$n$个不同输入；每个未见输入处的两个标签仍等可能，因此任意学习器
对随机目标的平均期望风险至少为$1/4$。于是存在一个固定的可实现目标，使算法的
平均期望风险至少为$1/4$。由$R\leq1$，有
$\mathbb E R\leq1/8+(7/8)\mathbb P\{R>1/8\}$，故失败概率至少为$1/7$。
这对每个$n$都成立，否定了分布无关的可实现PAC保证。
\end{proof}

最后一段证明并未使用标签数有限，因此有限Natarajan维对任意标签空间都是可学习性的
必要条件。有限标签条件用于相反方向的预测计数。上面的ERM样本量是一个充分条件，
并不声称它给出了可实现或不可知学习的最优样本复杂度。

\subsection{无限标签下的局限}
\label{sec:multiclass-erm-limitations}

当$Y$无限时，前面的$q=\binom K2$不再有限，预测计数不能沿原样使用。
这只是证明方法失效，还不足以说明Natarajan维本身失效。
要否定它的充分性，必须找到Natarajan维有限但不可学习的假设空间。

一个二维例子可以先展示Natarajan维遗漏了什么。考虑两个输入上的六种预测
\[
 (a,u),\ (b,u),\ (b,v),\ (c,v),\ (c,w),\ (a,w),
\]
其中$a,b,c$互异，$u,v,w$互异。固定任意一个预测，都能只改第一坐标得到另一个，
也能只改第二坐标得到另一个。例如$(a,u)$有邻居$(b,u)$和$(a,w)$。
但任何两个第一坐标标签至多共享一个第二坐标标签，因此找不到两对标签构成全部
四种组合。这个预测集合的Natarajan维为1，却在两个坐标上都持续保留了局部歧义。
单个有限例子本身仍然可学习；真正的反例需要任意高维的这种结构。

把二维例子推广到任意维，需要下面的组合构造。引理中将$V_m$的每个向量视为
定义在坐标集合$\{1,\ldots,m\}$上的预测函数，因此可以讨论它的Natarajan维。

\begin{lemma}[低Natarajan维的局部邻接构造]
\label{lem:low-natarajan-pseudocubes}
对每个整数$m\geq1$，存在有限标签集$Y_m$及有限非空集合$V_m\subseteq Y_m^m$，使
\begin{enumerate}
 \item $d_N(V_m)=1$；
 \item 对任意$v\in V_m$和$i\in\{1,\ldots,m\}$，存在$u\in V_m$满足
 \[
 u_i\neq v_i,\qquad u_j=v_j\quad(j\neq i).
 \]
\end{enumerate}
\end{lemma}

该引理引用Brukhim等的组合拓扑构造\scite{brukhim2022multiclass}本章不展开其构造证明。
第一条限制预先固定标签对的二元选择，第二条却保证每个预测在每个坐标上都有竞争者。
下面以这两条性质为依据，构造一个固定的不可学习假设空间。

\begin{theorem}[无限标签下有限Natarajan维并不充分]
\label{thm:natarajan-infinite-counterexample}
存在可数输入空间、可数标签空间和假设空间$\mathcal H$，使$d_N(\mathcal H)=1$，
但$\mathcal H$不是可实现PAC可学习的。
\end{theorem}

\begin{proof}
取引理~\ref{lem:low-natarajan-pseudocubes}中的$V_m$，令输入空间为互不相交的有限块
$X=\bigcup_{m\geq1}C_m$，其中$C_m=\{(m,1),\ldots,(m,m)\}$。
对每个局部标签$y\in Y_m$，用$(m,y)$表示带有块编号的标签。
例如，$(2,a)$和$(3,a)$是两个不同类别，即使它们都包含字母$a$。
还需要一个单独的类别，表示“输入不在这个假设所选的块内”。把这个类别记为$\star$，
并约定它不等于任何$(m,y)$。它是一个普通的预测标签，不表示未知标签、缺失值或任意取值。
完整标签空间为
\[
 Y=\{\star\}\cup\bigcup_{m\geq1}(\{m\}\times Y_m).
\]
对$v\in V_m$，定义
\[
 h_{m,v}(r,i)=
 \begin{cases}
 (m,v_i),&r=m,\\
 \star,&r\neq m.
 \end{cases}
\]
因此，$h_{2,v}$在块$C_2$内使用$v$给出的局部标签，在$C_1,C_3$等所有其他块上
都输出同一个“块外”类别$\star$。这样每个局部向量都被扩展成整个输入空间上的假设。
令$\mathcal H$由所有这些$h_{m,v}$组成。

两个来自不同块的输入不能被Natarajan打散：每个标签对至少有一个非$\star$标签，
但没有假设能同时在两个块上输出非$\star$标签。同一块内的两个输入也不能被打散：
若标签对包含$\star$，一处输出$\star$就迫使该块所有位置输出$\star$，不能独立选择；
若两对都不包含$\star$，则会要求$V_m$打散两个坐标，与其Natarajan维为1矛盾。
某个输入又能取局部标签或$\star$，所以$d_N(\mathcal H)=1$。

给定任意样本量$n$，取$m=2n$，令输入均匀分布于$C_m$，目标向量从$V_m$均匀选择。
对于样本中未出现的坐标$i$，即使额外告知学习器其余全部坐标的真实标签，
第$i$个坐标仍至少有两个等可能的候选标签：这些候选对应$V_m$中与其余坐标相同的
向量，每个向量的先验权重相同，且邻居条件保证候选数至少为2。
所以学习器在这个坐标上的平均错误概率至少为$1/2$。

至少一半输入未在样本中出现，故对随机目标、样本及学习器随机性的平均期望风险
至少为$1/4$。对目标取平均意味着存在一个固定$v\in V_m$也达到这个下界。
与前一定理相同，期望风险超过$1/8$的概率至少为$1/7$。
每个$n$都有这样的可实现分布，因而不存在分布无关的PAC样本量保证。
\end{proof}

障碍不在于标签无限就无法学习，而在于有限Natarajan维没有排除所有困难的预测结构。
接下来需要分别刻画两件事：整个假设空间是否一致收敛，以及是否存在能够学习它的算法。

\section{无限标签空间}
\label{sec:general-multiclass-dimensions}
\label{sec:multiclass-sample-complexity}

有限标签时，Natarajan维通过生长函数把一致收敛与可学习性联系起来。
第~\ref{sec:multiclass-erm-limitations}节说明，去掉标签数有限这一条件后，有限Natarajan维
已经不足以保证学习成功。此时还要区分两个问题：所有假设的经验风险能否同时准确
估计各自的期望风险，以及是否存在一个算法能够选出期望风险足够小的假设。
前者要求整个假设空间都受到控制，后者只要求算法作出合适的选择。

这两个问题需要不同的维度。一致收敛关注预测相对于真实标签是否正确，因而可以通过
图维转化为二值损失函数的VC维问题。可学习性则还要考虑标签本身携带的信息，
DS维通过预测之间只改变一个输入的关系，刻画有限样本无法排除的预测歧义。
下面先证明图维与一致收敛的关系，再说明为何可学习性需要单独分析。

\subsection{一致收敛}

期望风险只关心预测是否等于真实标签。即使错误标签有无限种，零一损失仍只有0和1。
因此可以为每个输入固定一个参考标签，再考察假设能否独立决定在哪些输入上与它相同。
这正是图维记录的结构。

\begin{definition}[图打散]
\label{def:graph-dimension}
若有限集$C\subseteq X$存在$f:C\to Y$，使每个$B\subseteq C$都对应某个
$h_B\in\mathcal H$，满足
\[
 h_B(x)=f(x)\ (x\in B),\qquad h_B(x)\neq f(x)\ (x\in C\setminus B),
\]
则称$C$被$\mathcal H$图打散。
\end{definition}

图打散描述一个给定点集上的预测能力。与VC维和Natarajan维一样，把可被打散的点集
大小取上确界，就得到整个假设空间的维度。

\begin{definition}[图维]
\label{def:graph-dimension-value}
设$\mathcal C_G$为所有被$\mathcal H$图打散的有限输入集组成的集合，定义
\[
 d_G(\mathcal H)=\sup_{C\in\mathcal C_G}|C|.
\]
\end{definition}

图打散只固定一个参考标签；不等于它时可以使用任何标签，而且可以随$B$改变。
Natarajan打散则事先固定替代标签。因此，Natarajan打散总能给出图打散，反过来不一定。

图~\ref{fig:graph-shattering-concrete}给出一个完整的小假设空间：输入只有
$x_1,x_2,x_3$，假设只有表中的四个。固定参考标签$a,u$后，前两个输入上“相同或不同”
的四种结果都能出现，所以这两个输入被图打散。第三个输入始终预测$z$，无法在参考
标签与其他标签之间切换，故这个假设空间的图维恰好为2。这里$x_1$偏离$a$时既可取$b$，
也可取$c$，说明图打散不要求预先固定唯一的替代标签。

\begin{figure}[htbp]
 \centering
 % Four hypotheses on a three-point domain: graph dimension exactly two.
\begin{tikzpicture}[x=1cm,y=1cm,
 cell/.style={draw=BookRule,minimum width=12mm,minimum height=6mm,font=\small},
 every node/.style={font=\small}]
 \node[anchor=west,font=\heiti\small] at (0,0.7) {参考预测：$f(x_1)=a,\ f(x_2)=u$};
 \foreach \x/\t in {0/假设,1.7/$x_1$,3.3/$x_2$,4.9/$x_3$,7.5/与参考预测的关系}{
  \node at (\x,0) {\t};
 }
 \foreach \r/\h/\a/\b/\txt in {1/$h_1$/a/u/两处相同,2/$h_2$/b/u/仅第二处相同,3/$h_3$/a/v/仅第一处相同,4/$h_4$/c/w/两处不同}{
  \node at (0,-0.75*\r) {\h};
  \node[cell] at (1.7,-0.75*\r) {$\a$};
  \node[cell] at (3.3,-0.75*\r) {$\b$};
  \node[cell,fill=BookPaper] at (4.9,-0.75*\r) {$z$};
  \node at (7.5,-0.75*\r) {\txt};
 }
 \draw[BookTeal,thick] (0.95,0.3) rectangle (4.05,-3.4);
 \node[align=center,text=BookTeal] at (2.5,-3.9) {$x_1,x_2$被图打散};
 \node[align=center,text=BookMuted] at (6.8,-3.9) {$x_3$始终预测$z$\\三个点不能被图打散};
\end{tikzpicture}

 \caption{图打散与图维的具体例子。绿色框内两列相对于参考预测实现全部四种相同或不同
 的组合；第三列固定不变。因此，在这个三点输入空间上，四个假设组成的空间图维为2。}
 \label{fig:graph-shattering-concrete}
\end{figure}
\FloatBarrier

为把多分类损失化为二分类函数，令
$\mathcal L_{\mathcal H}=\{(x,y)\mapsto\mathbb I\{h(x)\neq y\}:h\in\mathcal H\}$。
它的定义域是输入与标签的二元组，而输出只有0或1。

\begin{theorem}[图维与一致收敛]
\label{thm:graph-uniform-convergence}
对任意标签空间，$\operatorname{VC}(\mathcal L_{\mathcal H})=d_G(\mathcal H)$。
因此，零一损失的分布无关一致收敛当且仅当$d_G(\mathcal H)<\infty$成立。
\end{theorem}

\begin{proof}
若$C$被以$f$为参考图打散，则在点集$\{(x,f(x)):x\in C\}$上，
损失函数能够实现所有0/1预测，故损失函数空间的VC维至少为图维。

反过来，若一组$(x_i,y_i)$被损失函数空间打散，它不能包含同一个输入对应两个
不同标签的点，否则无法让两个损失同时为0。于是所有$x_i$互异；令$f(x_i)=y_i$，
损失取0与1分别对应预测等于和不等于$f(x_i)$，所以这些输入被图打散。
两种维度相等。由前章的VC一致收敛定理，二值损失函数空间满足分布无关的
一致收敛，当且仅当它的VC维有限，因此得到所述等价关系。
\end{proof}

图维解决了经验风险能否同时准确估计所有假设期望风险的问题。在这个保证下，
ERM选择经验风险最小的假设，就能得到接近假设空间最优期望风险的预测结果。
下面单独给出这一推论及其误差比较过程。

\begin{theorem}[一致收敛推出可学习性]
\label{thm:multiclass-uniform-pac}
设$\mathcal H$在零一损失下满足分布无关的一致收敛。对任意$\varepsilon,\delta\in(0,1)$，
若样本量$n$足以使任意联合分布$\mathcal P$下的样本$S\sim\mathcal P^n$满足
\[
 \mathbb P\!\left\{
 \sup_{h\in\mathcal H}|\hat R_S(h)-R_{\mathcal P}(h)|\leq\varepsilon/2
 \right\}\geq1-\delta,
\]
则任意ERM学习算法$A$均以至少$1-\delta$的概率满足
\[
 R_{\mathcal P}(A(S))\leq\inf_{h\in\mathcal H}R_{\mathcal P}(h)+\varepsilon.
\]
因此，$\mathcal H$不可知PAC可学习，也可实现PAC可学习。
\end{theorem}

\begin{proof}
考虑所有假设的经验风险与期望风险之差均不超过$\varepsilon/2$的样本事件。
它的概率至少为$1-\delta$，且其中的估计保证也适用于依赖样本选出的$A(S)$。
对任意$h\in\mathcal H$，有
\begin{align*}
 R_{\mathcal P}(A(S))
 &\leq\hat R_S(A(S))+\varepsilon/2\\
 &\leq\hat R_S(h)+\varepsilon/2\\
 &\leq R_{\mathcal P}(h)+\varepsilon.
\end{align*}
第一步和第三步分别使用$A(S)$与$h$的估计保证；第二步使用ERM的经验风险最小性。
对$h$取下确界就得到不可知PAC的期望风险保证。
由于一致收敛所需的样本量不依赖具体的$\mathcal P$，这个学习保证也是分布无关的。
可实现情形下，假设空间的最小期望风险为零，结论进一步成为
$R_{\mathcal P}(A(S))\leq\varepsilon$。
\end{proof}

结合这两个定理，有限图维足以保证可学习。但在无限标签下，有限图维不再是可学习的
必要条件。

一致收敛足以保证ERM学习，但它要求同时控制所有假设。以下例子中，标签本身能够
帮助算法识别目标假设，因而即使整个假设空间不一致收敛，仍然可以学习。

\label{ex:multiclass-cantor}

令$X=\mathbb N$。每个自然数的有限子集都作为一个标签，另设一个与所有集合标签不同的
类别$\star$。在本例中，它表示“输入不属于目标集合”，同样不是未知或缺失标签。
对每个有限$A$定义
\begin{equation}
 h_A(x)=\begin{cases}A,&x\in A,\\\star,&x\notin A,\end{cases}
 \qquad\mathcal H_C=\{h_A:A\subset\mathbb N\text{有限}\}.
 \label{eq:cantor-class}
\end{equation}
例如，标签$\{2,5\}$本身就告知哪些输入属于集合。因此，一旦观察到非$\star$标签，
学习器便能恢复整个目标假设；如果样本标签全为$\star$，就输出$h_{\varnothing}$。
这是一种利用标签中信息的数学构造，区别于一般语义类别。

\begin{theorem}[可学习性不推出一致收敛]
\label{thm:multiclass-erm-separation}
$\mathcal H_C$满足$d_N=1$及$d_G=\infty$。
它可实现PAC可学习，但不满足一致收敛，且存在不具有PAC保证的ERM选择规则。
\end{theorem}

\begin{proof}
设目标为$h_A$，$p=\mathcal P_X(A)$。上述算法见到一次非$\star$标签便具有零期望风险；
否则期望风险为$p$。若$p>\varepsilon$，后一事件的概率为
$(1-p)^n\leq e^{-n\varepsilon}$，故$n\geq\log(1/\delta)/\varepsilon$足够学习。

任意有限$C$都能以全$\star$为参考被图打散：要恰好在$B\subseteq C$上输出
$\star$，选择$A=C\setminus B$即可。因此图维无限，不能一致收敛。
另一方面，只要某个输入的预测是非$\star$标签，这个标签就唯一确定整个假设。
若两个输入被Natarajan打散，其中一个输入的两个标签至少有一个不是$\star$。
固定这个非$\star$标签后，另一个输入便无法再自由选择两个标签，产生矛盾。
单个输入能够出现两种标签，因此$d_N=1$。

构造坏ERM：当大小为$n$的样本标签全为$\star$时，输出
$h_{A_S}$，其中$A_S=\{1,\ldots,2n\}\setminus\{x_1,\ldots,x_n\}$。
其他可实现样本根据非$\star$标签恢复目标；不可实现样本任取经验风险最小者。
在全$\star$样本上经验风险为零，但对目标$h_{\varnothing}$和在
$\{1,\ldots,2n\}$上均匀的输入分布，它至少分错$n$个未见输入，期望风险至少为$1/2$。
每个$n$都有这样的分布，所以该ERM没有分布无关PAC保证。
\end{proof}

这里的好算法和坏算法都能使可实现样本上的经验风险为零，差别在于如何利用标签信息。
这类分离说明，多分类中“存在学习器”不再能通过“所有ERM都有保证”替代
\scite{daniely2015multiclass}它与第~\ref{sec:multiclass-erm-limitations}节的反例作用不同：那里的假设空间没有PAC学习器，
这里则只有某些ERM选择规则失败。

\subsection{可学习性}

集合标签的例子说明，图维无限并不排除存在有效的统计学习方法；另一方面，
有限Natarajan维也不能保证可学习。要寻找充要条件，需要一种比固定标签对更宽泛、
又不强求所有假设一致收敛的结构。前面六种预测的例子提供了线索：即使各输入没有
固定的两个替代标签，知道其他输入的预测后，仍可能无法确定遗漏输入的预测。

\label{sec:ds-dimension}

为了精确表达这种歧义，我们要求：集合中的每个预测向量，在每一个坐标上都能找到
一个只改变该坐标的预测向量。这种有限预测集合称为伪立方体。

\begin{definition}[伪立方体]
\label{def:pseudocube}
有限非空集合$V\subseteq Y^m$称为$m$维伪立方体，如果对每个$v\in V$和
每个$i\in\{1,\ldots,m\}$，都存在$u\in V$满足
$u_i\neq v_i$且$u_j=v_j$对所有$j\neq i$成立。
\end{definition}

六种预测构成二维伪立方体。二分类的全部$2^m$种预测也是$m$维伪立方体，
但多分类不要求每个坐标始终只使用两个标签。“有限”是定义的一部分：
无限邻接集合可能具有不同性质，不能将这一要求删除。

\begin{definition}[DS维]
\label{def:ds-dimension}
若$m$个互异输入组成的有序点集$C$满足$\mathcal H|_C$包含一个$m$维伪立方体，
则称$C$被DS打散。DS维$d_{DS}(\mathcal H)$定义为可被DS打散的有限点集大小的上确界。
\end{definition}

图~\ref{fig:ds-six-cycle}把第~\ref{sec:multiclass-erm-limitations}节的六种预测画成六个顶点。实线只改变第一坐标，
虚线只改变第二坐标。每个顶点都具有两种线，所以这六个预测构成二维伪立方体。
若把输入空间限定为这两个点，DS维恰为2，而Natarajan维只有1：六边形中没有
两对标签对应的四种完整组合。这个例子把两种维度的差别直接显示出来。

\begin{figure}[htbp]
 \centering
 % A six-cycle: every vertex has a neighbor in each of two coordinates.
\begin{tikzpicture}[
 vertex/.style={draw=BookRule,fill=BookPaper,minimum width=16mm,minimum height=7mm,font=\small},
 first/.style={BookBlue,thick},second/.style={BookAmber,thick,dashed}]
 \node[vertex] (v1) at (0,1.8) {$(a,u)$};
 \node[vertex] (v2) at (1.9,0.9) {$(b,u)$};
 \node[vertex] (v3) at (1.9,-0.9) {$(b,v)$};
 \node[vertex] (v4) at (0,-1.8) {$(c,v)$};
 \node[vertex] (v5) at (-1.9,-0.9) {$(c,w)$};
 \node[vertex] (v6) at (-1.9,0.9) {$(a,w)$};
 \draw[first] (v1)--(v2); \draw[second] (v2)--(v3);
 \draw[first] (v3)--(v4); \draw[second] (v4)--(v5);
 \draw[first] (v5)--(v6); \draw[second] (v6)--(v1);
 \draw[first] (3.2,1.5)--(4,1.5);
 \node[anchor=west,font=\small] at (4.1,1.5) {只改$x_1$的预测};
 \draw[second] (3.2,0.8)--(4,0.8);
 \node[anchor=west,font=\small] at (4.1,0.8) {只改$x_2$的预测};
 \node[anchor=north west,align=left,font=\small,text width=38mm] at (3.2,0.35)
 {每个顶点都连接一条实线和一条虚线，因此DS维为2。\par\smallskip
 找不到两对标签的全部四种组合，因此Natarajan维为1。};
\end{tikzpicture}

 \caption{DS维为2、Natarajan维为1的六种预测。每个顶点是一种在两个输入上的完整预测，
 两种线型分别表示只改一个坐标的邻接。每个顶点在两个方向上都有邻居。}
 \label{fig:ds-six-cycle}
\end{figure}
\FloatBarrier

观察$m-1$个坐标的标签后，若遗漏坐标仍有邻居，就无法仅从其余标签确定它。
要求集合中每个向量、每个坐标都有邻居，排除了只在某个特殊向量周围制造少数分支的情况。
图~\ref{fig:learning-multiclass-dimensions}比较三种维度关注的对象；图中的DS部分
只展示邻接方式，并非一个完整伪立方体。

\begin{figure}[htbp]
 \centering
 % 多分类维度关注对象的层次：标签对、参考标签、局部邻接。
\begin{tikzpicture}[
  panel/.style={book node,minimum width=31mm,minimum height=42mm,inner sep=4pt},
  tag/.style={draw=BookRule,rounded corners=1pt,minimum width=6mm,minimum height=5mm,font=\scriptsize},
  note/.style={font=\kaishu\scriptsize,text=BookMuted,align=center}
]
  \node[panel,fill=FigureInput!7] (nat) {};
  \node[panel,fill=FigureModel!7,right=6mm of nat] (graph) {};
  \node[panel,fill=FigureLoss!7,right=6mm of graph] (ds) {};

  \node[font=\heiti\small,anchor=north] at ([yshift=-3mm]nat.north) {Natarajan维};
  \node[font=\heiti\small,anchor=north] at ([yshift=-3mm]graph.north) {图维};
  \node[font=\heiti\small,anchor=north] at ([yshift=-3mm]ds.north) {DS维};

  \foreach \i/\x in {1/-9,2/0,3/9}{
    \node[tag,fill=white] at ([xshift=\x mm,yshift=6mm]nat.center) {$a_{\i}$};
    \node[tag,fill=white] at ([xshift=\x mm,yshift=-4mm]nat.center) {$b_{\i}$};
    \draw[BookBlue,thick,<->] ([xshift=\x mm,yshift=3mm]nat.center) -- ([xshift=\x mm,yshift=-1mm]nat.center);
  }
  \node[note,anchor=south] at ([yshift=2mm]nat.south) {每个坐标预先固定\\一对可切换标签};

  \foreach \i/\x in {1/-9,2/0,3/9}{
    \node[tag,fill=white] at ([xshift=\x mm,yshift=6mm]graph.center) {$r_{\i}$};
    \node[tag,fill=white] at ([xshift=\x mm,yshift=-4mm]graph.center) {$\neq r_{\i}$};
    \draw[FigureModel,thick,->] ([xshift=\x mm,yshift=3mm]graph.center) -- ([xshift=\x mm,yshift=-1mm]graph.center);
  }
  \node[note,anchor=south] at ([yshift=2mm]graph.south) {固定参考标签；离开后\\具体标签可以变化};

  \node[tag,fill=white] (v1) at ([xshift=-8mm,yshift=5mm]ds.center) {$abc$};
  \node[tag,fill=white] (v2) at ([xshift=8mm,yshift=5mm]ds.center) {$dbc$};
  \node[tag,fill=white] (v3) at ([xshift=-8mm,yshift=-6mm]ds.center) {$aec$};
  \node[tag,fill=white] (v4) at ([xshift=8mm,yshift=-6mm]ds.center) {$abf$};
  \draw[FigureLoss,thick] (v1)--(v2);
  \draw[FigureLoss,thick] (v1)--(v3);
  \draw[FigureLoss,thick] (v1)--(v4);
  \node[note,anchor=south] at ([yshift=2mm]ds.south) {研究完整标记向量间\\只改一个坐标的邻接};
\end{tikzpicture}

 \caption{多分类维度的三种观察方式。Natarajan维固定每个坐标上的两个标签，图维固定
 参考标签，DS维考察预测向量间只改变一个坐标的邻接。DS列仅为局部邻接示意；
 伪立方体还要求其中每个向量在每个坐标上都有邻居。}
 \label{fig:learning-multiclass-dimensions}
\end{figure}

\begin{theorem}[多分类维度的比较]
\label{thm:multiclass-dimension-comparison}
对任意标签空间，有$d_N(\mathcal H)\leq d_{DS}(\mathcal H)\leq d_G(\mathcal H)$。
二分类时三者均等于VC维。若$|Y|=K<\infty$，则$d_N<\infty$也蕴含$d_G<\infty$。
\end{theorem}

\begin{proof}
Natarajan打散给出的全部二元选择构成伪立方体，因此第一道不等式成立。
对于$m$维伪立方体，固定一个向量$v$作为参考。给定希望改变的坐标集合，
逐个沿这些坐标移动到邻居；每一步只改变当前坐标，因此已改变的坐标保持不变，
未要求改变的坐标始终与$v$相同。最终向量恰好在指定坐标上不等于$v$，故图打散成立。

二分类中，任何坐标的替代标签唯一，三种结构都要求全部二元预测，所以等于VC维。
有限标签时，若图打散$m$个输入，至少需要$2^m$个不同预测向量。
由定理~\ref{thm:multiclass-growth-bound}，当$d_N=d<\infty$时，
$2^m\leq G_{d,K}(m)$。固定$d,K$后右侧是多项式，故$m$有有限上界。
\end{proof}

更精细地分析最后一个不等式，可得$d_G=O(d_N\log K)$；这里需要标签数有限。
对无限标签，第~\ref{sec:multiclass-erm-limitations}节的反例有$d_N=1$、$d_{DS}=\infty$，说明第一道不等式可以相差无穷。
前面集合标签的假设空间$\mathcal H_C$则满足$d_{DS}=1$：任何出现非$\star$标签的
预测向量都唯一确定假设，因而不可能在另一个坐标方向上拥有邻居；全$\star$向量独自
也不能构成伪立方体。结合单点可以打散可知DS维为1，而图维无限，故后一道不等式也不能反向替代。

伪立方体的局部邻居意味着未见坐标仍存在无法从其余标签排除的答案。
这一结构可以直接给出学习失败概率的下界。

\begin{theorem}[无限DS维的不可学习性]
\label{thm:infinite-ds-lower-bound}
若$d_{DS}(\mathcal H)=\infty$，则对每个样本量$n$及任意学习算法，存在
$\mathcal H$可实现分布，使输出假设的期望风险超过$1/8$的概率至少为$1/7$。
\end{theorem}

\begin{proof}
取$2n$个DS打散的输入和其上的有限伪立方体$V$。令输入均匀分布，并从$V$均匀
选择目标预测向量，每个向量选取一个实现它的假设。对未见坐标，即使知道其余全部
标签，也至少有两个等可能的候选答案，因此平均错误概率至少为$1/2$。
至少一半输入未被观察到，所以平均期望风险至少为$1/4$。
对目标取平均后选出达到该下界的固定目标，再用
$\mathbb E R\leq1/8+(7/8)\mathbb P\{R>1/8\}$，即得结论。
\end{proof}

有限DS维的充分性需要构造一个学习器，不能直接使用整个假设空间的一致收敛。
Brukhim等的构造先输出有限的标签列表，再从列表中选择单个标签；两步合起来得到
一种样本压缩方案。这里先明确引用的构造结论，再证明它如何给出期望风险保证
\scite{brukhim2022multiclass}

标签列表是一个映射$\mu$，它为每个输入$x$给出有限集合$\mu(x)\subseteq Y$。
例如，$\mu(x)=\{\text{正常},\text{广告}\}$表示暂时保留这两个答案。
列表覆盖样本，是指每个样本标签都属于相应列表；这还不是最终的分类器。
构造的第二步利用这些列表重构一个输出单个标签、且与原样本一致的假设。

\begin{lemma}[有限DS维的压缩构造]
\label{lem:finite-ds-compression}
设$d=d_{DS}(\mathcal H)<\infty$。存在通用常数$C>0$，使得对每个整数$n\geq1$，
所有大小不超过$n$的$\mathcal H$可实现样本均存在如下压缩表示：从样本中选取
一个允许重复、保留次序的带标签序列，长度至多
\[
 r_n=\min\!\left\{n,\left\lceil
 C(d+1)^{3/2}\log^2\!\bigl(2n(d+1)\bigr)\right\rceil\right\},
\]
并记录原样本的大小。对每个原样本大小，存在预先固定的重构映射，能够仅根据该压缩
序列重构一个与原样本一致的假设。重构假设不要求属于$\mathcal H$。
\end{lemma}

这是原论文定理36的一个放宽了常数和对数因子的表述。该构造的两项关键结果是：
有限DS维允许用少量样本重构覆盖原样本的标签列表；给定列表后，有限Natarajan维
允许再用少量样本重构一致的单标签预测。原论文引理39和40分别证明这两步。
拼接两段压缩序列即可先重构列表、再重构分类器，且$d_N\leq d$保证第二步也由$d$控制。
取原论文中的参数$t=\lceil\sqrt{d+1}\rceil$，所得长度不超过上述对数平方上界。
若这个上界超过$n$，直接保存整个样本即可。列表构造所需的超图定向论证引用原论文，
以下证明不将这一核心构造当作已在本章证明的结果。

有了压缩表示，可以沿用第~\ref{chap:binary-learnability}章的思路：可能的压缩序列数量有限，因此能够同时控制
所有重构假设的期望风险。这里的计数取决于压缩长度，不取决于标签空间的大小。

\begin{theorem}[有限DS维的期望风险保证]
\label{thm:finite-ds-learning-rates}
在所需可测性条件成立的设定下，设$d=d_{DS}(\mathcal H)<\infty$，$r=r_n<n$，并记
\[
 B_n(\delta)=\log\frac{2(n+1)(r+1)n^r}{\delta}.
\]
存在允许非恰当输出的学习算法，使对任意$\delta\in(0,1)$，以下保证成立：
\begin{enumerate}
 \item 对任意$\mathcal H$可实现分布$\mathcal P$，以至少$1-\delta$的概率，
 \[
 R_{\mathcal P}(A^{\mathrm{real}}(S))
 \leq\frac{B_n(\delta)}{n-r}.
 \]
 \item 对任意联合分布$\mathcal P$，以至少$1-\delta$的概率，
 \begin{align*}
 R_{\mathcal P}(A^{\mathrm{agn}}(S))
 &\leq\inf_{h\in\mathcal H}R_{\mathcal P}(h)+\frac rn\\
 &\quad+\sqrt{\frac{B_n(\delta)}{2(n-r)}}
       +\sqrt{\frac{\log(2/\delta)}{2n}}.
 \end{align*}
\end{enumerate}
\end{theorem}

\begin{proof}[证明（使用引理~\ref{lem:finite-ds-compression}）]
长度至多$r$的有序压缩序列最多有$(r+1)n^r$种索引选择。
再计入原样本大小的$n+1$种可能，重构方式至多有
$M=(n+1)(r+1)n^r$种。这里是在观察数据之前枚举索引，而不是预先枚举标签。

固定一种索引选择，并设它实际使用了$q\leq r$个不同样本。
给定这$q$个样本的值后，重构假设已经确定；其余$n-q$个样本仍独立同分布。
若重构假设的期望风险超过$u$，它在这些剩余样本上全部正确的条件概率至多为
$(1-u)^{n-q}\leq e^{-(n-r)u}$。
可实现学习器选择一个与全部样本一致的重构假设，压缩引理保证这样的选择存在。
对至多$M$种重构方式的失败概率求和，取$u=B_n(\delta)/(n-r)$，便得到第一项结论。
若$u\geq1$，结论直接由零一损失下的期望风险不超过1成立。

不可知学习器在这$M$种重构结果中选择经验风险最小者。固定一种重构方式，
记其在未参与重构的样本上的经验风险为$\hat R_{\mathrm{rest}}$。
由单侧Hoeffding不等式及上述条件独立性，以至少$1-\delta/2$的概率，
所有重构结果同时满足
\[
 R_{\mathcal P}(h)\leq\hat R_{\mathrm{rest}}(h)
       +\sqrt{\frac{\log(2M/\delta)}{2(n-r)}}.
\]
从全部样本中去掉$q$个观测，经验风险最多增加$q/n$：
将全部样本上的平均损失写成保留部分与删除部分的加权平均，即可得到
$\hat R_{\mathrm{rest}}(h)\leq\hat R_S(h)+q/n$。
于是上式可以用$\hat R_S(h)+r/n$代替$\hat R_{\mathrm{rest}}(h)$。

对任意预先固定的比较假设$g\in\mathcal H$，取样本中被$g$正确分类的全部观测。
这部分样本可实现，压缩引理保证某个重构假设在这些观测上全部正确，因而其经验风险
不超过$\hat R_S(g)$。若这部分为空，该不等式对任意假设都成立。
最小化重构结果的经验风险因此保证
$\hat R_S(A^{\mathrm{agn}}(S))\leq\hat R_S(g)$。
再对固定$g$使用Hoeffding不等式，以至少$1-\delta/2$的概率有
\[
 \hat R_S(g)\leq R_{\mathcal P}(g)
       +\sqrt{\frac{\log(2/\delta)}{2n}}.
\]
合并两项概率保证，再令固定比较假设的期望风险趋近于
$\inf_{g\in\mathcal H}R_{\mathcal P}(g)$，便得到第二项结论。
\end{proof}

对固定$d$，压缩长度$r_n$只按$\log^2 n$增长，因此$r_n/n\to0$。
上述两项期望风险上界都随样本量增加趋于相应的最优值。
忽略关于$n,d$的对数因子，可实现期望风险为
$\widetilde O((d^{3/2}+\log(1/\delta))/n)$，
不可知超额期望风险为
$\widetilde O(\sqrt{(d^{3/2}+\log(1/\delta))/n})$。
这些是统计保证；压缩表示的存在没有保证找到它的程序可计算或运行时间为多项式。

将这一构造与前面的必要条件合起来，可以得到与二分类学习基本定理相对应的结论。
二分类中由VC维统一刻画的性质，在一般多分类中分成了图维和DS维两组条件；
只有标签有限时，这两组条件的有限性才重新等价。

\begin{theorem}[多分类学习基本定理]
\label{thm:ds-characterization}
对任意标签空间，在所需可测性条件成立且允许非恰当学习时，以下命题等价：
\begin{enumerate}
 \item $d_{DS}(\mathcal H)<\infty$；
 \item $\mathcal H$可实现PAC可学习；
 \item $\mathcal H$不可知PAC可学习。
\end{enumerate}
此外，$d_G(\mathcal H)<\infty$当且仅当零一损失满足分布无关的一致收敛；
这两个等价条件均蕴含上述三条，但一般不存在反向蕴含。
\end{theorem}

\begin{proof}[证明（使用有限DS维的学习器构造结果）]
若DS维有限，定理~\ref{thm:finite-ds-learning-rates}分别给出可实现和不可知的分布无关期望风险保证，
选择足够大的样本量即可达到任意$\varepsilon,\delta$。
不可知PAC保证蕴含可实现PAC保证，而无限DS维由前一定理排除可实现PAC学习。
三个命题因而等价。图维与一致收敛的等价性见定理~\ref{thm:graph-uniform-convergence}，
反向蕴含的失败由定理~\ref{thm:multiclass-erm-separation}给出。
充分性中的学习器构造作为已引用的外部结果使用。
\end{proof}

因此，一般多分类中应分别使用两条结论：图维有限当且仅当零一损失一致收敛，
DS维有限当且仅当存在PAC学习器。前者推出后者，后者不推出前者。
若标签数有限，则$d_N$、$d_{DS}$、$d_G$的有限性、一致收敛以及上述两种PAC可学习性
全部等价；二分类时三种维度均退化为VC维。这些结论由定理~\ref{thm:multiclass-dimension-comparison}
和定理~\ref{thm:finite-label-learnability}得到。

定性刻画也不直接提供最优样本复杂度或高效实现，不能把二分类样本量公式中的VC维
简单替换为DS维。

\section{可计算PAC学习}
\label{sec:computable-multiclass-learning}

\subsection{统计保证与程序实现}

与第~\ref{sec:computable-binary-learning}节相同，本节依次区分统计可学习、程序可计算、
多项式时间可实现三个层次。变化的是标签从0和1扩展为$Y$中的元素，训练与预测仍然
必须停机，期望风险的成功标准也不变。

沿用定义~\ref{def:binary-computable-pac}，取$X=\mathbb N$，并为标签约定可计算编码和
可判定的相等关系。多分类不可知CPAC学习器是一个对每份有限样本停机的程序，
输出一个对每个输入都能停机返回$Y$中标签的程序，并满足
\[
 \mathbb P_{S\sim\mathcal P^n}\!\left\{
 R_{\mathcal P}(A(S))\leq\inf_{h\in\mathcal H}R_{\mathcal P}(h)+\varepsilon
 \right\}\geq1-\delta
\]
对所有允许分布及所有足够大的$n$成立。要求输出属于$\mathcal H$时称为恰当学习；
允许输出其他可计算假设时称为非恰当学习。样本量上界还要求可计算的strong CPAC，
与前章使用同一约定，不表示运行时间一定为多项式。

\begin{theorem}[多分类可计算ERM的充分条件]
\label{thm:multiclass-computable-erm}
若$\mathcal H$存在对所有有限样本都停机、输出假设可有效求值的恰当ERM，并且
$d_G(\mathcal H)<\infty$，则$\mathcal H$恰当地不可知CPAC可学习，且有可计算的样本量上界。
有限标签时，可以将有限图维条件替换为有限Natarajan维。
\end{theorem}

\begin{proof}
定理~\ref{thm:graph-uniform-convergence}给出损失函数空间的有限VC维，因此有显式的
一致收敛样本量上界。给定精度$\varepsilon/2$，ERM的两次偏差比较给出超额期望风险
至多为$\varepsilon$。训练和预测的可计算性由题设保证。
固定假设空间的有限维度上界可作为常数写入样本量程序。有限标签时，有限Natarajan维
由定理~\ref{thm:multiclass-dimension-comparison}推出有限图维。
\end{proof}

这与前章可计算ERM的定理对应：二分类使用VC维，多分类使用损失空间的图维，
有限标签下也可使用Natarajan维。只有统计维度有限而没有可计算ERM时，仍不能
据此断言存在这样的程序。

\subsection{可计算Natarajan维}

有限Natarajan维保证：给定足够多的输入及每个输入处的两个标签，总有一种二元选择
不能被实现。要利用这一条件编写程序，需要能够在有限时间内找出这样的选择。

\begin{definition}[可计算Natarajan维]
\label{def:computable-natarajan-dimension}
若存在总可计算程序$W$，对任意$d+1$个互异输入$x_i$及标签对$(a_i,b_i)$、
$a_i\neq b_i$，都输出$s\in\{0,1\}^{d+1}$，使得
\[
 \nexists h\in\mathcal H\quad
 \forall i:\ h(x_i)=\begin{cases}a_i,&s_i=0,\\b_i,&s_i=1,\end{cases}
\]
则称$d$为可计算Natarajan维的一个上界。最小的这种非负整数记为
$c_N(\mathcal H)$；不存在时记为$\infty$。
\end{definition}

普通Natarajan维只要求缺失的预测存在，可计算版本还要求一个统一程序找出它。
所以$d_N\leq c_N$；当标签只有两种时，这一概念回到前章的有效VC维。

证明的充分方向使用Gourdeau等的有效扩展引理：若$c_N(\mathcal H)=d<\infty$且
标签有限，可构造包含$\mathcal H$的假设空间$\mathcal H'$，使
$d_N(\mathcal H')\leq d+1$，并且有限输入上的全部可实现预测及其实现假设可有效求得
\scite{gourdeau2025multiclass}以下明确区分这一外部构造与随后的统计推导。

\begin{theorem}[有限标签的非恰当CPAC刻画]
\label{thm:computable-multiclass-characterization}
在上述有效编码模型中，若$2\leq|Y|<\infty$，则$\mathcal H$非恰当地不可知
CPAC可学习，当且仅当$c_N(\mathcal H)<\infty$。
\end{theorem}
\begin{proof}[证明（充分性使用有效扩展引理）]
必要性可通过可计算版本的没有免费的午餐论证得到。设存在CPAC学习器$A$，
固定一个使精度$1/16$、失败概率$1/16$的PAC保证成立的样本量$n$。
给定$2n$个输入及每处两个不同标签，枚举全部二元选择作为候选目标；
对每个候选目标，枚举来自这$2n$个输入的所有长度$n$的有序样本，运行$A$，
再计算输出在这$2n$个输入上的错误比例。所有集合有限，训练和预测都停机，
因而每个候选目标下的平均期望风险可计算。

随机二元目标的未见输入论证保证，某个候选目标的平均期望风险至少为$1/4$。
程序找出这样一个目标并返回它的选择向量。该目标不可能属于$\mathcal H$的限制，
否则PAC保证会使平均期望风险至多为$1/16+1/16<1/4$。
这就给出维度至多为$2n-1$的有效见证。这里$n$作为固定常数写入程序，
不要求从任意学习器描述中自动算出它。

充分性使用上述扩展$\mathcal H'$。枚举有限样本上的可实现预测并选择错误最少的一种，
再输出其可计算实现，得到一个可计算ERM。由于标签有限且$d_N(\mathcal H')$有限，
本章有限标签定理保证它不可知PAC学习$\mathcal H'$。
又因$\mathcal H\subseteq\mathcal H'$，有
$\inf_{h\in\mathcal H'}R_{\mathcal P}(h)\leq\inf_{h\in\mathcal H}R_{\mathcal P}(h)$，
故同一算法也满足相对于$\mathcal H$的不可知保证，但其输出未必属于$\mathcal H$。
\end{proof}

此处学习器取确定性程序，概率来自随机样本；若采用额外随机化的计算模型，必须同时
约定随机位的使用方式，不能直接假定可以枚举所有随机运行过程。

上述必要性证明没有使用标签集合有限，因此无限标签下CPAC学习仍要求$c_N<\infty$。
充分性中的有效扩展则用到了有限标签的预测枚举，不能直接推广。
更强的是，本章不可学习反例的构造并不因只检查标签对就消失：一个有限的计算维度
条件是否充分，需要单独证明，不能从统计定义的相似性推出。

有限DS维解决的是统计可学习性，不自动提供伪立方体搜索、列表学习和单标签归约的
可执行程序。因此不能把DS刻画直接改称为多分类CPAC定理。
对于具体假设空间，一条可验证的充分路线仍然是：给出可计算ERM，并证明损失空间
一致收敛；有限标签和有限Natarajan维可以保证后一个条件，任意标签下有限图维也可以。

\subsection{高效学习与优化误差}

可计算程序仍然可能运行很久。与前章一样，高效PAC学习还要控制样本量、训练时间、
预测时间及输出表示长度。多分类需额外注意标签编码：标签数为$K$时，类别编号的长度
约为$\log K$；若输出标签本身是集合或其他结构，其描述长度也必须计入成本。

近似ERM的统计分析完全保留。若算法输出$A(S)\in\mathcal H$，并满足
\[
 \hat R_S(A(S))\leq\inf_{h\in\mathcal H}\hat R_S(h)+\alpha,
\]
且所有假设的经验风险与期望风险之差不超过$\eta$，则
\[
 R_{\mathcal P}(A(S))\leq\inf_{h\in\mathcal H}R_{\mathcal P}(h)+2\eta+\alpha.
\]
证明就是定理~\ref{thm:approximate-erm-risk}的三步比较，不依赖标签有几个。
其中$2\eta$来自抽样估计，$\alpha$来自优化误差。若能以多项式时间求得
$\alpha\leq\varepsilon/2$的解，并用多项式数量的样本使$\eta\leq\varepsilon/4$，
便得到高效不可知PAC学习的一条充分路线。

\subsection{从枚举到有效算法}

前章用布尔合取说明，逐个枚举假设很慢，并不意味着所有算法都很慢。
多分类也可以用一个明确例子看出这一区别。令输入空间为$X=\{1,\ldots,p\}$，
标签空间为$Y=\{1,\ldots,K\}$，并允许全部分类函数$\mathcal H=Y^X$。
共有$K^p$个假设，但经验风险最小化不必枚举它们。

对每个输入$x$，统计样本中它与各标签共同出现的次数
$N_{x,y}=|\{i:x_i=x,\ y_i=y\}|$。在$x$处选择出现次数最多的标签，
平局时选择编号最小者；没有出现过的输入也按这一约定预测1。
不同输入处的选择相互独立，所以可以分别求解。

\begin{theorem}[有限输入上的多分类表学习]
\label{thm:multiclass-table-erm}
上述计数算法是$Y^X$上的恰当ERM，需要$O(n+pK)$次计数和比较操作，
输出一张含$p$个标签的预测表。对任意$\varepsilon,\delta\in(0,1)$，若
\[
 n\geq\frac{2}{\varepsilon^2}
       \left(p\log K+\log\frac2\delta\right),
\]
则以至少$1-\delta$的概率，其期望风险不超过$\inf_{h\in\mathcal H}R_{\mathcal P}(h)+\varepsilon$。
\end{theorem}

\begin{proof}
任意假设$h$分对的样本总数为$\sum_{x=1}^pN_{x,h(x)}$。
每个$x$独立选取使$N_{x,y}$最大的$y$，就使这个和最大，因此经验风险最小。
初始化$pK$个计数器、扫描$n$个观测、逐行比较，得到所列操作数。
又因$|\mathcal H|=K^p$，有限假设空间的Hoeffding分析使
$\sup_h|R_{\mathcal P}(h)-\hat R_S(h)|\leq\varepsilon/2$的失败概率至多为
$2K^p e^{-n\varepsilon^2/2}$。令其不超过$\delta$，再使用ERM比较，即得结论。
\end{proof}

这里计算代价按$p$和$K$衡量；计数器的位数至多为$O(\log(n+1))$，因此按比特计费
也仍为这些参数的多项式。但若把$K$看成由一个很短程序描述的巨大标签空间，
就不能把$O(K)$误称为标签编码长度的多项式。输出预测表的长度为$O(p\log K)$，
也计入学习资源。

\section{本章小结}
\label{sec:multiclass-learning-summary}

把标签编码成若干二元位，可以把预测任务拆成二分类子问题，但并未消除对这些子问题
复杂度、错误累积与输出一致性的分析。有限标签下，Natarajan维限制有限输入上的不同
预测数目，使二分类中的计数、对称化与ERM比较仍然适用。因此，在本章的可测性条件下，
有限Natarajan维、分布无关PAC可学习性与一致收敛三者等价。

标签空间无限时，需要把两个问题分开。图维刻画零一损失函数空间的VC维，决定能否用
同一份样本同时准确评价全部假设。DS维则刻画统计可学习性，允许学习算法只利用足够的
预测信息，而不必准确评价整个假设空间。有限Natarajan维此时只保留必要性，不能独自
保证学习成功。有限图维推出可学习性，反方向不成立；这正是二分类基本定理不能原样
推广的原因。

这一区分也影响算法选择。经验风险同样达到最小值的学习器，可能具有不同的泛化表现，
所以仅写出ERM目标不足以解释所有多分类学习现象。进一步要求训练与预测程序停机时，
还要引入有效的组合条件；有限标签下的可计算Natarajan维提供了相应刻画，却不自动保证
多项式运行时间。下一章将把评价对象扩展为一般损失，并从算法稳定性继续分析
一致收敛之外的学习保证。

\paragraph{延伸阅读}
Daniely等关于多分类可学习性与ERM的研究，适合用来进一步理解“存在成功学习器”与
“任意ERM都能成功”的区别。可重点对照文中的不同ERM表现与本章图维的作用，检查
算法在多个经验最优假设之间怎样选择，会如何影响样本量\scite{daniely2015multiclass}

Brukhim等给出了DS维对多分类可学习性的刻画。本章已经展开列表预测、压缩与单标签预测
之间的统计推导；若要补全所引用的组合构造，应继续阅读原文中的伪立方体与列表学习分析。
这些构造是深层的组合结果，不能由Natarajan维的有限标签证明直接替代\scite{brukhim2022multiclass}

计算方向可阅读Gourdeau、Lechner与Urner的工作，重点核对标签编码、输出假设的求值要求
以及恰当与非恰当学习的区别。它提供了与第~\ref{sec:computable-binary-learning}节
有效VC维分析相对应的多分类视角\scite{gourdeau2025multiclass}
